\clearpage
\fchapter{El sistema de mando programado}
Controladores de procesos (PLC, VARIADOR, PANTALLAS, ETC). Aquí se aglutina la
documentación generada sobre el sistema de control programado.
Para cada uno de ellos se realizarán los apartados siguientes:
\begin{enumerate}
	\item Funciones asignadas. Se indicará: Descripción y características del controlador.
	\item Diagramas de control o Grafcet adaptados al controlador seleccionado. En este sentido se trata
	      de la realización de los diagramas de control del apartado anterior. Estos diagramas se han
	      realizado en la fase del diseño. A partir de los mismos se han generado las instrucciones que
	      componen el programa de control y las rutinas de verificación.
	\item Listado de variables. Donde los comentarios adquieren especial relevancia, para ser buscados
	      e interpretados.
	\item Listado del Programa o programas de control. Se ofrecerán estructurados en bloques de
	      programación. Incluso se organizarán los subprogramas.
	\item Listados de programas que componen las rutinas de verificación (OB30, OB80, etc), útiles en
	      operaciones de puesta en marcha y todo tipo de mantenimiento. En este tipo se incluyen las
	      rutinas de detección de condiciones de error, y transmisión de mensajes a las unidades interfaz
	      hombre-máquina.
	\item Listado de plantillas de datos, para la verificación de funciones. Por ejemplo, configuración
\end{enumerate}
